%% ==========================================================================
%% Anhang & Anhangsverzeichnis (Exakt gemäß IU-Richtlinien)
%% ==========================================================================

% --------------------------------------------------------------------------
% Anhangsverzeichnis (verpflichtend bei mehreren Anhängen)
% Laut IU-Richtlinien: Unmittelbar vor dem eigentlichen Anhang eingefügt,
% enthält die Bezeichnungen der Anhänge ohne Seitenzahlen.
% Erscheint als 'Anhangsverzeichnis' im Hauptinhaltsverzeichnis.
% --------------------------------------------------------------------------
\phantomsection
\section*{Anhangsverzeichnis}
\addcontentsline{toc}{section}{Anhangsverzeichnis}

\vspace{0.8cm}

\noindent
\textbf{Anhang A:} Interviewleitfaden und Erhebungsinstrumente\\[0.4cm]
\textbf{Anhang B:} Ergänzende statistische Auswertungen und Datensätze\\[0.4cm]
\textbf{Anhang C:} Quellcode und technische Konfigurationen

\newpage

% --------------------------------------------------------------------------
% Anhang A
% --------------------------------------------------------------------------
\phantomsection
\section*{Anhang A: Interviewleitfaden und Erhebungsinstrumente}
\label{anhang:A}

Dieser Anhang enthält den für die empirische Erhebung verwendeten Interviewleitfaden. Im Textteil der Bachelorarbeit muss mindestens einmal auf diesen Anhang verwiesen werden (z.\,B. \enquote{siehe Anhang~A}).

\lipsum[1-2]

\newpage

% --------------------------------------------------------------------------
% Anhang B
% --------------------------------------------------------------------------
\phantomsection
\section*{Anhang B: Ergänzende statistische Auswertungen und Datensätze}
\label{anhang:B}

Hier können umfangreiche Tabellen, Fragebögen oder statistische Auswertungen dokumentiert werden, die den fortlaufenden Lesefluss im Hauptteil stören würden.

\lipsum[3]
\lipsum[1-2]